\documentclass[conference]{IEEEtran}

\usepackage[T1]{fontenc}
\usepackage[utf8]{inputenc}
\usepackage{cite}
\usepackage{amsmath,amssymb,amsfonts}
\usepackage{graphicx}
\usepackage{booktabs}
\usepackage{hyperref}
\hypersetup{hidelinks}

\title{Galaxy Graph Lab: IEEE Report Skeleton}
\author{
\IEEEauthorblockN{Santiago Florido Gomez}
\IEEEauthorblockA{Galaxy Graph Lab\\
Paris, France}
}

\begin{document}
necesit
\maketitle

\begin{abstract}
This document defines the IEEE-style report structure for Galaxy Graph Lab.
It is intended as the starting point for documenting the graph model,
constraint formulation, implementation decisions, and validation results.
\end{abstract}

\begin{IEEEkeywords}
graph theory, constraint solving, puzzle modeling, pygame
\end{IEEEkeywords}

\section{Introduction}

Galaxy Graph Lab studies graph-based representations for structured puzzles and
constraint-driven search. This report follows the standard IEEE conference
layout so the project can grow into a technical paper without restructuring the
document later.

The current implementation baseline includes a Python package, a Pygame entry
point, and a reporting workspace for documenting the mathematical formulation
and the experimental results. The software stack is described together with the
report skeleton so that implementation and documentation evolve in the same
repository.

\section{Mathematical Model}

In order to define the problem to be modeled in a clear and mathematically precise way, it is first necessary to establish the rules of the game and describe the dynamics of each instance. This provides the basis for subsequently formulating the constraints that enforce those rules and give rise to the optimization system solved for each instance.

\section{Methodology}

This section should explain how the mathematical model is mapped to code. It is
the right place to document:

\begin{itemize}
\item the solver strategy,
\item the graph data structures,
\item the Pygame interface responsibilities, and
\item the validation workflow used during development.
\end{itemize}

\section{Results}

This section should summarize the first validation runs for Galaxy Graph Lab.
Typical IEEE-ready content includes benchmark instances, execution time, solver
success rate, and representative visual outputs generated by the interface.

\section{Conclusion}

The report now uses an IEEE structure that is ready for iterative expansion.
Future revisions should replace the placeholders with the final formulation,
implementation details, and quantitative evaluation of the project.


\begin{thebibliography}{00}
\bibitem{pygame}
Pygame Community, ``Pygame documentation.'' [Online]. Available:
\url{https://www.pygame.org/docs/}
\end{thebibliography}

\end{document}
